Die Ergebnisse wurden schon ausreichend in der Auswertung vorgestellt und müssen nicht nochmal zusammengefasst werden. 
\paragraph{Wellenlänge des Lasers}
Die Messung scheint sehr genau zu sein, die Abweichung von \(0.4\sigma\) zum Literaturwert zeigt gute Übereinstimmung mit dem Herstellerwert. Potenzielle Fehlerquellen sind das nicht synchrone Ablesen von Motorskala (Feinuhr) und Diskriminator, sowie eine ungenaue Einstellung der Nullposition der Messuhr. 

\paragraph{Brechungsindex}
Auch hier liegt die Abweichung auf einem sehr guten Niveau, bei \(0.9\sigma\). Dies liegt auch daran, dass die Fitergebnisse der Ausgleichsgeraden alle identisch sind und dementsprechend kein statistischer Fehler einfließt. 

\paragraph{Kohärenzlänge}
Große Schwierigkeiten hat während dem Versuch das Aufnehmen bzw. erst einmal Finden der Interferenzbilder auf dem Oszilloskop bereitet. Es war ziemlich schwierig, das \enquote{Wellenpaket} aus \cref{fig:ledinterferogramm} zu erwischen. Des weiteren habe ich erst in der Auswertung wirklich verstanden, was wir dort überhaupt gemessen haben. 

Für die Kohärenzlänge der LED existiert kein Literaturwert. Die Größenordnung von wenigen Mikrometern ist jedoch typisch für LEDs. Externe Quellen geben für grüne LEDs eine Kohärenzlänge von etwa \(\pm 4\,\mu\mathrm{m}\)\footnote{https://www.opto-e.com/de/resources/tutorials/coherence-artifacts-on-imaging-systems} an, was konsistent mit dem gemessen Wert und der Größenordnung übereinstimmt. Die Messung ist daher konsistent und plausibel.
